% !TeX root = ../kazymyr-thesis.tex

\chapter{Approach}\label{chapter:approach}

This chapter describes the implementation details added in the course of this
project. It covers all the components that were added to extend the received
simulator so that it supports the HYBRID mode.

\section{Starting point and contributions}\label{sec:approach-delta}

The base simulator described in Section~\ref{sec:prelim-simulator} is
existing work~\cite{zacouris2025m2ts,zacouris2026mxs}. Two tables make clear
which parts are inherited and which are new, before the design is described.
Table~\ref{tab:approach-received} lists the components that this thesis
builds on. Each of them is used as it was received.
Table~\ref{tab:approach-added} in the next section lists what is new. None of
its entries has a counterpart in the base simulator.

\begin{table}[htbp]
  \centering
  \captionsetup{format=plain,justification=centering}
  \caption{Taken from the base simulator.}
  \label{tab:approach-received}
  \footnotesize
  \begin{tabular*}{\textwidth}{@{\extracolsep{\fill}}>{\raggedright\arraybackslash}p{0.26\textwidth} >{\raggedright\arraybackslash}p{0.61\textwidth}@{}}
    \toprule
    Execution engine & Fixed time step, IO and CPU processor pools \\
    \midrule
    ACID runtime & Lock table, conservative 2PL with cautious waiting \\
    \midrule
    BASE runtime & Dirty set, its (model, resource) key, fan-out
    collapse, and the reconciler pool that clears it \\
    \midrule
    Rule expansion & Bookkeeper transactions and multi-model CPR \\
    \midrule
    Workload generator & Seeded transactions under the \texttt{EQUILIBRIUM} profile \\
    \bottomrule
  \end{tabular*}
\end{table}

The baseline configuration of Section~\ref{sec:exp-setup} follows the
reference setup of the base simulator for the system and workload
parameters, but not for every value. The runs are five times longer,
steady-state metrics ignore an initial warm-up window, and the two oracle
thresholds have no counterpart in the reference setup at all. Each of these
differences is stated with the baseline table and justified there. Keeping
the remaining values at the reference setup is what makes the validation in
Section~\ref{sec:x1-validation} meaningful.




\section{Introduction of Hybrid Mode}\label{sec:approach-why-hybrid}

The purpose of the hybrid mode is to turn inconsistency into a threshold that
an operator can set. Simulator suppports two semantcs. First is ACID, the slow but
consistent approach. Its dirty fraction $d_k$ is always zero, but its throughput
drops sharply as the number of resources and systems that one ACID
transaction has to reach grows~\cite{zacouris2026mxs}. BASE is the other
side. It lets the work that has to be propagated to other models be
distributed to reconciliation processors, which clear the updates
eventually. In BASE, however, the inconsistency of the system cannot be
controlled at all, and it can grow so fast that the reconciliation capacity
is unable to keep up. Our approach combines the two modes and adds a
mechanism for inconsistency control. The operator states a critical
dirty threshold $c$, and the system spends as much ACID execution as it needs
to stay under that level and runs BASE the rest of the time.

Table~\ref{tab:approach-added} lists the components of the hybrid mode that
were introduced in this thesis, with a reference to the section where each of
them is described in more detail.

\begin{table}[htbp]
  \centering
  \captionsetup{format=plain,justification=centering}
  \caption{Added in this thesis.}
  \label{tab:approach-added}
  \footnotesize
  \begin{tabular*}{\textwidth}{@{\extracolsep{\fill}}>{\raggedright\arraybackslash}p{0.28\textwidth} >{\raggedright\arraybackslash}p{0.59\textwidth}@{}}
    \toprule
    HYBRID execution mode (Section~\ref{sec:approach-dispatchers}) & The mode
    becomes a per-transaction tag instead of a run-wide configuration \\
    \midrule
    Shared-resource arbiter (Section~\ref{sec:approach-arbiter}) & Lookup
    that connects the ACID and BASE bookkeeping on shared resources, so that
    neither mode overwrites the work of the other \\
    \midrule
    Claim table (Section~\ref{sec:approach-claim-table}) & Records the
    locks used by BASE parent transactions \\
    \midrule
    Fixed-share oracle
    (Section~\ref{sec:approach-fixed-share-oracle}) & Tags ACID with a
    configured probability $p_{\mathrm{acid}}$ \\
    \midrule
    Adaptive oracle (Section~\ref{sec:approach-adaptive-oracle}) & Switches
    to ACID when the dirty fraction reaches the threshold, and back to BASE
    once the inconsistency has recovered \\
    \midrule
    Hybrid metrics & Mode switches, ACID share, blocking and ACID visit
    duration \\
    \bottomrule
  \end{tabular*}
\end{table}


\section{Received simulator allows only one mode per run}\label{sec:approach-pipeline}

The simulator can run ACID or BASE, but never both. A run binds its
configured concurrency model through the \emph{runtime assembler}, and there
is one assembler per mode. It builds every mode-dependent collaborator at the
start of the simulation, so all transactions run in the one mode that was
assembled at the beginning. Figure~\ref{fig:pipeline} shows how the
mode-specific components are wired. While explaining how the two modes are
started, this section also follows the path that one ACID or one BASE
transaction takes, because the hybrid mode reuses that path.

At the start of the simulation the runtime assembler (Step 1 in
Figure~\ref{fig:pipeline}) creates a single wiring, either ACID or BASE. Only
one path exists in the simulator at a time. For example, when the simulator
is configured in ACID mode, none of the BASE components are built, and vice
versa. The assembler is the only place that both modes share.

Admission (Step 2) has two parts, and the rest of this chapter keeps them apart.
The first part is the same in both modes. The admission controller takes the
next transaction from the pending queue, but only while the number of
transactions in flight is below the capacity limit. The second part is the
\emph{mode admission step}, and it differs between the modes.

BASE has no locking step and admits the transaction straight into execution.
ACID first acquires the whole lock set of a parent transaction at once, or waits
while holding none of it. A transaction counts as \emph{admitted} only once it
has passed the mode admission step. A new transaction that arrives from the
workload is not admitted yet. It only joins the pending queue. The concurrency
strategy and the transaction executor (Steps 3 and 4) then drive the phases
\texttt{LOCK}, \texttt{IO}, \texttt{CPU} and \texttt{COMPLETE}, and BASE skips
\texttt{LOCK}.\label{sec:approach-execution}

In BASE mode the processor pool is split into two sets. The first is the
foreground processors, which handle the initial parent transactions. The second
is the reconciler pool, which propagates the changes to the other models. The
simulator allows both the number of foreground IO and CPU processors and the
number of reconciler processors $R$ to be configured. In BASE, $R$ is subtracted
from the pool of foreground processors. In ACID the value of $R$ is ignored.
After a parent transaction is processed, rule expansion (Step 5) decides what
happens next.\label{sec:approach-completion} ACID parks the parent as waiting
for its group and spawns its bookkeepers, which have to go through the same
process of acquiring locks. BASE commits the parent immediately and then starts
its bookkeepers independently, without waiting for them.

\begin{figure}[htbp]
  \centering
  \includegraphics[width=\textwidth]{fig-received-pipeline.pdf}
  \captionsetup{format=plain,justification=centering}
  \caption{The received pipeline.}
  \label{fig:pipeline}
\end{figure}

One property of the ACID mode is important. The unit of locking is the group and
not the transaction, because the root preclaims its own set at admission while
its bookkeepers extend that set afterwards.\label{sec:approach-group-locking}
Two groups can therefore wait on each other, and the received simulator resolves
this by aborting one of them.

All five components exist in an ACID form and a BASE form, and the form is fixed
at assembly. A run therefore inherits one execution mode and keeps it to the
end.\label{sec:approach-seams}

The idea of the hybrid mode is to keep both modes and to wrap each of these
components in a dispatcher that holds one collaborator of each kind.

\section{Implementing the hybrid mode}\label{sec:approach-hybrid}

A third execution mode, \texttt{HYBRID}, is added next to ACID and BASE. It
is not a third set of semantics but a configuration in which both received
sets are live at the same time. What the hybrid mode adds is the choice of
which mode a transaction runs in.

\paragraph{Dispatchers wrap both paths.}\label{sec:approach-dispatchers} A third
runtime assembler builds both mode-specific sets in full and wraps each pair.
Every component of Section~\ref{sec:approach-pipeline} is then filled by a
dispatcher that holds one ACID collaborator and one BASE collaborator, and
forwards each call to the one that the transaction names. Both ACID and BASE
paths are left unchanged. The two admission controllers share the scheduler, the
processors, the I/O model, the queue manager and the capacity limit, and the two
bookkeeper coordinators share the dirty set and the lock table. The dispatchers
and the arbitration checks are built by the hybrid assembler only, so a pure
ACID or pure BASE run still executes the received code as in received simulator.

\paragraph{Introduction of tag and mode
  oracle.}\label{sec:approach-oracle-interface}\label{sec:approach-tag}\label{sec:approach-tagging-point}
Since both modes are wired into the simulator at the same time, the decision
which path to take has to move out of the assembler. Our approach introduces a
tag for each transaction. Only two tags exist: \emph{\{ACID, BASE\}}. Each transaction
receives its tag once, and the tag is propagated to the bookkeepers of that
transaction. Choosing an execution mode per transaction at runtime is not new in
itself. Hybrid Transactional Replication~\cite{kobus2018htr} and
PolyCert~\cite{couceiro2011polycert} already do it, but there all modes give the
same consistency, so the choice is never about controling the inconsistency levels
(Section~\ref{sec:rw-mode-selection}).

The tag is chosen by a component called the \emph{mode oracle}. It is
implemented as a Java interface with a single method, \texttt{assignMode()}.
This thesis introduces two mode oracles, which are described in
Section~\ref{sec:approach-oracles}. The oracle is called between the two parts
of admission. Once the admission controller has taken an untagged parent
transaction from the pending queue, the mode oracle returns the mode in which
the transaction should run, and the dispatcher hands it to the \emph{mode
  admission step}, either ACID or BASE.

Tags are immutable. They cannot be withdrawn or changed. Once a transaction has
received its tag, it stays in that mode for its whole life, and so do all of its
bookkeepers. The same holds when a transaction is aborted. It loses all its
progress, but not its tag.

\paragraph{Sharing one resource pool, and the limits of the naive
  design.}\label{sec:approach-resource-layer}\label{sec:approach-mixing-hazard}
One set of resources is created at the beginning and is never divided again.
At the baseline that is $5000$ resources, divided equally between the $10$
models. An ACID-tagged transaction may lock any of the $5000$
resources, and a BASE-tagged transaction may dirty any of the same $5000$.

Simply wrapping both modes is a naive design here. It separates the path a
transaction takes, but the resources that transaction needs are still shared. The
problem is that the lock table used by ACID and the dirty set used by BASE
do not communicate. As Figure~\ref{fig:hybrid-resources} shows, both modes
ignore the progress of the other when they work on the same resource. ACID
locks are visible only to other ACID transactions, and BASE dirty marks only
to other BASE transactions. This would break the isolation that ACID
promises, because a locked resource could be overwritten while it is locked.

\begin{figure}[H]
  \centering
  \includegraphics[width=\textwidth]{fig-hybrid-resources.pdf}
  \captionsetup{format=plain,justification=centering}
  \caption{Resource management in the hybrid mode.}
  \label{fig:hybrid-resources}
\end{figure}

A resource that is ACID-locked and BASE-dirty at the same time is in an
undefined state, because each half relies on an assumption that the other
half breaks. Sharing the pool is therefore only half of what a hybrid run
needs. The other half is a shared-resource arbiter. It connects the
bookkeeping of the two pure modes and defines what a mode does when a
resource it needs is in use by the other mode. The next section addresses this problem.

\section{The shared-resource arbiter}\label{sec:approach-arbiter}

Section~\ref{sec:approach-mixing-hazard} left the hybrid mode with one open
problem. Both halves write to the same resources and ignore the work of the
other half. The lock table is asked only whether a lock is free. The dirty
set is asked only whether a repair is in progress. Nothing in that
arrangement stops a resource from being locked and dirty at the same time.

This section describes the shared-resource arbiter, the mechanism
introduced to solve this isolation problem between the ACID and BASE
bookkeeping.

\paragraph{Arbiter implementation.}\label{sec:approach-arbiter-lookup} The
arbiter is not a single class that holds references to all resources or
locks. It reuses two existing records, the lock table and the dirty set,
and adds a third one, the claim table.
Figure~\ref{fig:arbiter-layer} shows how the three records are connected.
Every lookup starts from the same place, the name of a resource in the
shared pool (column 1). An ACID transaction, a BASE parent and a BASE repair
all start from that name when they want to use resource $\mathit{res}_k$.
The name is then mapped to a key (column 2). The lock table and the claim
table both use the lock id, which the granularity mapping computes from the
resource. A run may be configured with fewer locks than resources, so one
lock id can cover several resources. The dirty set uses the dirty key, the
pair of the owning model and the resource, so it always has one entry per
resource. Under its key each record keeps one entry (column 3). The lock
table lists the ACID transactions holding the lock, the claim table lists
the BASE parents writing under it, and the dirty set holds the
repair that is still outstanding on the resource. Because both keys are
computed from the same resource name, any path can look up any record
without extra bookkeeping.

\begin{figure}[H]
  \centering
  \includegraphics[width=\textwidth]{fig-arbiter-layer.pdf}
  \captionsetup{format=plain,justification=centering}
  \caption{The shared lookup the arbiter is built on.}
  \label{fig:arbiter-layer}
\end{figure}

The colours in Figure~\ref{fig:arbiter-layer} show who writes and who reads
each record. The lock table is written and read by ACID as before, and is
now also read by both BASE paths. The claim table is written by BASE parents
and the dirty set by BASE repairs, and both are read by ACID. Each record is
therefore still written by one side only. What the arbiter adds are reads
across the two modes.

These reads are placed in the mode admission step (Step 2 in
Figure~\ref{fig:pipeline}), which is extended to also check the state of the
opposite mode. For example, when an ACID transaction goes through its mode
admission step, the arbiter checks whether any resource it needs is dirty or
claimed, by looking it up in the dirty set and in the claim table. A
transaction that fails this check is not admitted. It has already left the
pending queue, but it goes to the waiter registry or to the deferred parent
queue instead of into execution. BASE repairs are the one
exception, because they are created later, at rule expansion (Step 5). The
idea is that every path checks a resource for work of the other mode before
it uses it.

The components introduced by the arbiter, both to resolve the lookup and to
handle a mode that finds a resource in use by the other mode, are the
following.

\begin{bullets}
  \item[Claim table.]\label{sec:approach-claim-table} One entry per lock
  under which an admitted BASE parent is currently writing. A parent enters
  its claims when it is admitted and removes them when it completes.

  \item[Waiter registry.] The ACID transactions parked on a resource they
  need that is currently dirty or claimed. No locks are acquired and no state
  is changed.

  \item[Deferred parent queue.] The BASE parents that met an ACID lock at
  admission.

  \item[Deferred repair queue.] The BASE repairs whose target was
  ACID-locked when the rule that created them fired.
\end{bullets}

Only the claim table is a new record. The other three components \emph{Waiter
  registry}, \emph{Deferred parent queue}, \emph{Deferred repair queue} hold the
transactions that found a resource in use by the other mode, until that resource
is free again.

A large simplification is important to mention. The resources are split
between models, so in a real system a lookup takes time, because each model
that holds the resource has to be reached. These network calls may be slow
and can sometimes fail. The simulator sees all resources and models as one
entity, so a lookup has no latency at all and always succeeds. Modelling the
real system more closely would add a lot of overhead. The simulator
therefore implements the mechanism in an ideal setting, and it would be less
performant in an actual multi-model system.

\paragraph{The invariant.}\label{sec:approach-arbiter-rule} The main goal of
the hybrid implementation, and of the arbiter itself, is to enforce one
rule. A resource is never held by ACID and written by BASE at the same time.
This keeps the guarantees of ACID transactions intact and addresses the idea
of safe adaptation between ACID and BASE posed by RQ1. Because a resource
always has one clearly defined state, no work is lost or overwritten.
Section~\ref{sec:approach-deadlock} covers this in more detail.

\paragraph{Shared arbiter lookup position.}\label{sec:approach-enforcement}
To enforce the invariant, the pipelines of both ACID and BASE are extended
with a lookup for work of the opposite mode. An ACID transaction is tested
before it takes any lock. A BASE parent is tested at admission, before it
claims anything. A BASE repair is tested at rule expansion, before its
target is marked dirty.

\paragraph{ACID extended admission step.} In Figure~\ref{fig:hybrid-queueing-acid}
we can see how the intial ACID path was extended to perform look up for BASE work.  \todo{add note that this disagrm is taken from
  "Modeling ACID and BASE Semantics for Cross-Model Consistency-Preserving
  Updates" and extended here} An ACID transaction is tested against the claim table and the dirty
set together, and both answers have the same effect. If any resource it needs is
dirty or claimed, the transaction is paused immediately and added as a waiter to
the waiter registry. It never acquires any locks, which is an important step in
the later deadlock-freedom argument. The transaction stays in the
\texttt{PENDING} state, but it is removed from the pending queue.
Figure~\ref{fig:hybrid-queueing-acid} shows that it enters the waiter registry
rather than the lock queue, so it holds nothing while it waits.

Two events make a resource available again. A repair that finishes clears
the dirty entry, and a BASE parent that completes drops its claims. A waiter
on a dirty resource is registered under the resource, and a waiter on a
claim under the lock id, because claims are kept per lock. All waiters
registered under the cleared key are then notified and retried. A root parked at
admission returns to the pending queue for a fresh attempt, and a bookkeeper
parked in mid-execution goes through the existing retry path.

\begin{figure}[htbp]
  \centering
  \includegraphics[width=\textwidth]{fig-hybrid-queueing-acid.pdf}
  \captionsetup{format=plain,justification=centering}
  \caption{ACID meeting a resource it may not touch.}
  \label{fig:hybrid-queueing-acid}
  \medskip
  \begin{figurekeys}
    \item[1] The transaction is registered as a waiter and taken out of the
    pending queue. The test runs before the lock phase, so it holds no lock
    while it waits.
    \item[2] Once the resources are clean again the registry notifies every
    transaction waiting on them. Each one returns to the pending queue,
    where the test is repeated.
  \end{figurekeys}
\end{figure}

\paragraph{BASE extended admission step.} The two BASE paths are tested
separately, because they reach a resource at different times.
Figure~\ref{fig:hybrid-queueing-base-parent} shows the parent path, and
Figure~\ref{fig:hybrid-queueing-base-bookkeeper} shows the repair path.
The parent is the foreground work and is tested at mode admission. Every resource
it needs is checked before it claims any of them. If one of them is
ACID-locked, the parent claims nothing and goes on the deferred parent
queue. Inside this queue a BASE parent does not wait for a specific ACID
transaction. It is retried on every drain, so if the resource is still
locked, even by a lock acquired in the meantime, the parent is deferred
again rather than overrunning the lock.

\begin{figure}[htbp]
  \centering
  \includegraphics[width=\textwidth]{fig-hybrid-queueing-base-parent.pdf}
  \captionsetup{format=plain,justification=centering}
  \caption{A BASE parent meeting an ACID lock at admission.}
  \label{fig:hybrid-queueing-base-parent}
  \medskip
  \begin{figurekeys}
    \item[1] A parent whose resources are all free returns to the pending
    queue and is asked the same question at the next admission step.
    \item[2] On completion the claims are dropped and every transaction
    waiting on those resources is notified.
    \item[3] The queue is retried on every admission. A parent whose resource
    is still locked stays deferred and goes to the back of the queue.
  \end{figurekeys}
\end{figure}

The repair is tested just before it would mark its target dirty. If the
target is ACID-locked, the repair is not created at all. It goes on the
deferred repair queue, the resource stays clean, and the transaction that
triggered it commits. Deferred does not mean waiting. The repair is retried
systematically, but there is no guarantee of when it can actually proceed.
This avoids a dependency of a bookkeeper on an ACID transaction, which could
otherwise lead to a deadlock, and it breaks no promise of eventual
consistency. When the repair is finally created, it
is created in full and processed.

The claim table is read only by ACID. A BASE parent never checks it, and a
dirty resource does not stop it either, so BASE work never waits for other
BASE work on the same resource. This follows the idea of Salt
Isolation~\cite{xie2014salt}, where BASE transactions may see each other's
unfinished work, while ACID transactions are kept apart from it
(Section~\ref{sec:rw-static}).

\begin{figure}[htbp]
  \centering
  \includegraphics[width=\textwidth]{fig-hybrid-queueing-base-bookkeeper.pdf}
  \captionsetup{format=plain,justification=centering}
  \caption{A BASE repair meeting an ACID lock at rule expansion.}
  \label{fig:hybrid-queueing-base-bookkeeper}
  \medskip
  \begin{figurekeys}
    \item[1] Once the lock is gone the target is marked dirty and the repair
    reconciliation starts.
    \item[2] On every retry a repair whose target is still locked is kept in
    the queue.
  \end{figurekeys}
\end{figure}


Both deferred queues are polled rather than woken. Nothing signals an entry
when its lock disappears. Parents are retried once per admission step, before
anything is admitted. Repairs are retried after every batch of repairs, after
every ACID group completion, and once per simulated tick. The pass per tick
is the one that matters, because a lock is also released by a group that
aborts, and an abort is not a completion. Whatever released the lock, the
entry is retried within one tick.

\paragraph{Hybrid mode limits the parent BASE
  transactions.}\label{sec:approach-base-cost} Making a BASE parent wait is
where BASE loses performance, and this cannot be avoided. BASE is meant to be
the fast mode, which processes the transaction on its own model almost
immediately and resolves the consequences in the background. In the hybrid
mode a BASE parent can be held back at admission, so it loses part of that
speed.

A BASE parent waits only at admission, only for as long as an ACID lock is
held on one of its resources, and only while holding nothing itself. It
never blocks after it has been admitted, it never waits on another BASE
transaction, and it never waits for reconciliation. It therefore stays much
faster than an ACID transaction, but slower than in a pure BASE run.

\paragraph{Where the arbiter is deliberately
  over-cautious.}\label{sec:approach-conservatism} Two imprecisions are
built into the simulator by design.

\begin{bullets}
  \item[C1. Lock buckets hide free resources.] The two BASE checks and the
  ACID claim check read the lock of a resource and not the resource itself.
  Where a run is configured with fewer locks than resources, one lock
  covers a bucket of them, so BASE work is set aside whenever any
  resource in that bucket is locked. Transactions therefore wait in cases
  where they could have proceeded.

  \item[C2. Starvation of deferred work.] A deferred entry whose target
  stays locked can stay on its queue indefinitely. No entry is
  forgotten, but there is no ageing rule that eventually gives a
  long-deferred entry priority over the ACID transaction keeping its
  resource busy, so under a resource that is hot for a long time
  that work can starve. This is accepted rather than solved. Eventual
  consistency places no bound on when a repair runs.
\end{bullets}


\section{Mode oracles}\label{sec:approach-oracles}

Section~\ref{sec:approach-oracle-interface} gave a brief introduction of what an
oracle is. This section provides the overview of the two that were built during
the course of the thesis. They share the interface, the arbiter, the tagging
rules and the dispatchers, and they differ only in one aspect, how decision
about tag is made.

\paragraph{Short intodcution to both oracles.} During this thesis work two type
of oracles were implemented. The both have the full hybrid mechanism underneath,
same arbiter, same tagging and same dispatchers. The difference we focused on
that one is static and other is dynamic. The static one is
\texttt{FixedShareModeOracle} it is configured once at the start of the
simulation and doesn't rely on any more internal or external information. On the
other hand we have \texttt{AdaptiveModeOracle}, it focuses on leveragin the
system's incosistency levels $d_k$ sampled at transaction admission to make a
decision which tag the transaction should receive.

\paragraph{The fixed-share oracle in
  details.}\label{sec:approach-fixed-share-oracle} The fixed-share oracle returns
ACID with a configured probability $p_{\mathrm{acid}}$ and BASE otherwise
$p_{\mathrm{base}} = 1 - p_{\mathrm{acid}}$, from a seeded generator, so that a
run can be reproduced. The draw is random but the share of ACID/BASE
transactions is not.

It was first hybrid design that we focused on and for a practical reason. Since
$p_{\mathrm{acid}}$ can be set from the simulation start, it is easy to
configure and fixed-share oracle worked as proof of concept for hybrid mode.
Apaprt from that, we planed to use the this type of oracle in the evaluation, to
answer the RQ3 and compare if adaptive configuartion can outperform static one.
It was intented as the ablation baseline.

One design property of this oracle sets how it can be compared with the adaptive
one. What $p_{\mathrm{acid}}$ fixes is the share of \emph{transactions} tagged
to be in ACID mode. It fixes nothing about the share of \emph{time} the system
spends in ACID, because the mode of a transaction decides how long that
transaction occupies the system.

In short, this fixed-share oracle generates a controlled split of ACID and BASE
tags assigned to pending incoming transactions.

\paragraph{The adaptive oracle in details.}\label{sec:approach-adaptive-oracle} The
adaptive oracle reads the live dirty fraction $d_k$ at the instant of the decision
and switches between two modes. It starts always in BASE mode, since we want to
priorities the throughput whenever possible. As as soon the systems level of
inconsistency goes over the
critical dirty fraction threshold $c$ that was introdcued as part of this
thesis, adaptvie oracle swithches to the ACID mode, and all next transactions
would be tagged as ACID.
To get back to base, the level on incosistency $d_k$ has to drop to the dirty
recovery threshold $r$.

\begin{align}
  \text{BASE} \rightarrow \text{ACID} \quad & \text{when} \quad d_k \ge c
  \label{eq:enter-acid}                                                   \\
  \text{ACID} \rightarrow \text{BASE} \quad & \text{when} \quad d_k \le r
  \label{eq:leave-acid}
\end{align}

with the interval $[r, c]$ acting as the hysteresis band
(Figure~\ref{fig:oracle-fsm}).

The oracle locks one mode for the whole system until the switch conditions are
met. Every new transaction admission is tagged with whatever that mode is at that
instant. Work that is already in flight keeps the mode it was admitted under, by
the lifetime rules of Section~\ref{sec:approach-tag}. A mode change is therefore
not an event that the running transactions experience.

Incosistency level $d_k$ is the only metric the oracle uses, so every mode
transition has a single clear cause. The Figure~\ref{fig:oracle-fsm}
demonstrates how the adaptive oracle behaves, when the mode switch happens and
what tag each transaction will receive at some simuatlion point.

Past $c$ no further BASE work is
admitted and only the BASE transactions already in flight still add dirt,
and past $r$ only the reconciliations already running carry on. The two
edges are not symmetric in how far they run. The overshoot past $c$ is
small, but the $r$ undershoot can be large, because the oracle reads $d_k$
only at admission (Section~\ref{sec:approach-sampling}) and admissions are
rare in ACID, so reconciliation keeps draining while nothing is looking. More
details you can find in Section~\ref{sec:recovery-mechanism}.

\begin{figure}[htbp]
  \centering
  \includegraphics[width=\textwidth]{fig-hysteresis-band.pdf}
  \captionsetup{format=plain,justification=centering}
  \caption{The hysteresis band against a dirty-fraction trace.}
  \label{fig:oracle-fsm}
  \medskip
  \begin{figurekeys}
    \item[1] The overshoot past $c$ is small. It is carried by the BASE
    transactions that were already in flight when the mode changed,
    because every arrival after the switch is tagged ACID.
    \item[2] The undershoot past $r$ can be large. The oracle reads $d_k$
    only at admission (Section~\ref{sec:approach-sampling}) and ACID
    admissions are rare, so reconciliation keeps draining while nothing
    is looking (Section~\ref{sec:recovery-mechanism}).
  \end{figurekeys}
  \vspace{-12pt}
\end{figure}

\paragraph{The critical dirty and recovery
  thresholds.}\label{sec:approach-parameters-oracle} The oracle has two
parameters. The critical dirty or also incosistency level threshold $c$ is the
fraction of resources that allowed to be dirty before the oracle stops issuing
BASE tags. It is the value of simulator configuration about how much
inconsistency the deployment is willing to accept.

The recovery threshold $r$ is the fraction the system must reconcile back down
to before BASE resumes. This reconciliation happens because oracle switches to
ACID mode, which is slower and fully consistent semantic. That allows to avoid
additional dirtiness and gives reconcile processors time to catch up. Both of
the threshold are fractions of the resources $N$, and $r$ is configured directly
rather than derived from $c$, subject to $0 \le r \le c$.

The two extremes of that range are worth naming, because they are policies and
not just values. At $r = c$ the band has zero width and the only anti-flap
mechanism in the design is gone (Section~\ref{sec:approach-nofilter}). The
simulator accepts this configuration, but it is not a recommended setting and no
experiment in this thesis uses it. At $r = 0$ the rule says reconcile everything
before resuming BASE, which differs from a small $r$, because an ACID visit then
ends only when the dirty set is empty, so its length is set by the reconciler
pool alone.

\paragraph{The oracle reads incosistency level $d_k$ at the
  admission.}\label{sec:approach-sampling}\label{sec:approach-samplers} The oracle
reads the dirty fraction inside the call that assigns a mode, that is, once per
admission of an untagged root transaction, nothing else triggers a reading. The
simulation tick does not forces it, and no reading is taken while the
system is progressing but admits nothing.

This ties the rate at which the oracle sees the system dirty fraction to the rate at which
transactions are admitted. Under \texttt{EQUILIBRIUM} the transaction population is
held constant, so admissions tied to completions and the reading rate of the
oracle is approximately the throughput. That means when adaptive oracle is in
ACID mode the throughput drops significantly, and with that the rate of
inconsistency $d_k$ lookup and Section~\ref{sec:effective-recovery}
measures consequenses of that.

It is important being explicit that this is not how the reported measurements of
dirty resources fraction are taken. The reading of the oracle is part of the
system under study. The reported measurements come from a separate reader that
runs once at the end of each simulated tick regardless of mode or load.

Reading $d_k$ only at admission is still enough, because a reading matters only when a
transaction has to be tagged. A reading once per tick would change the mode of
the oracle earlier, but it would not affect any transaction earlier. Under
\texttt{EQUILIBRIUM} a new root transaction is created only when an in-flight
one finishes. In ACID mode these transactions are long, and they hold the
admission capacity while reconciliation clears the dirty set in the background.
The first BASE transaction is therefore admitted at the same moment with either
kind of reading, and it sees the same $d_k$.

\paragraph{The adaptive oracle is studied only under a stationary
  load.}\label{sec:approach-stationary-load} The adaptive oracle is designed for
a load that changes during a run. When more work arrives, the dirty fraction
grows faster and the oracle moves more work to ACID. When the load falls, it
returns to BASE. All experiments in Chapter~\ref{chapter:experiments} use the
\texttt{EQUILIBRIUM} profile, which keeps the number of in-flight transactions
constant (Section~\ref{sec:prelim-workload-profiles}). The load is therefore
stationary, and the oracle has no change in load to react to. It still switches
modes, but only in response to the dirty fraction that its own decisions
produce. The received simulator also has a profile with a varying load,
\texttt{TERMINAL}, but it was still under development while the experiments
were run, so it is not used.

This has two consequences. First, the comparison with the fixed-share oracle for
RQ3 (Section~\ref{sec:x3-frontier}) is made in a setting that is not ideal for
the adaptive oracle, because a fixed ACID share that is tuned for one
configuration never has to change under a constant load. Second, every run uses
one fixed setup, for example one reconciler count $R$. The results show that the
oracle keeps the dirty fraction close to critical incosistency threshold $c$ in
each of these setups. They do not show how the oracle behaves when the load
changes during a run. Testing the oracle under a varying load is left for future
work (Section~\ref{sec:future-work}).

\paragraph{Collision probability of ACID
  transactions and dirty resources.}\label{sec:approach-blocking-cost} The
adaptive oracle introduces the critical dirty threshold $c$. This means that
only about a fraction $c$ of the resources can be inconsistent at the same time.
Knowing that value, we can model a useful quantity, the probability that one
ACID transaction has to wait on a dirty resource.

One parameter has to be explained before the derivation. The baseline value
$L = 50$ (\texttt{maxTransactionResources}) is the maximum number of resources
one transaction can touch. The simulator draws the resource count $\ell$ of each
transaction uniformly from $1$ to $L$. The derivation assumes that these
resources are independent of each other and that each of them is dirty with
probability $d$.

In this setup a single resource is clean with probability $1-d$. A transaction
with $\ell$ resources is clean only if every one of them is clean, so the
probabilities multiply.

\begin{equation}
  P_{\text{clean}}(d) \;=\; \underbrace{(1-d)\cdot(1-d)\cdots(1-d)}_{\ell \text{ resources}}
  \;=\; (1-d)^{\ell}
  \label{eq:clean-probability}
\end{equation}

Blocking is the complementary event, in which at least one of the $\ell$
resources is dirty.

\begin{equation}
  P_{\text{block}}(d) \;=\; 1 - P_{\text{clean}}(d) \;=\; 1-(1-d)^{\ell}
  \label{eq:blocking-probability}
\end{equation}

Figure~\ref{fig:blocking-probability} plots this probability against the dirty
fraction $d$ from $0$ to $0.12$ for two cases. The first is the worst case
$\ell = L = 50$. The second is the average over all transaction sizes $\ell$
from $1$ to $50$. The figure shows that blocking is not a rare event but the
normal case for an ACID transaction admitted while the dirty set is large. The
system is not always at the worst state $d_k = 0.10$. But even at the baseline
recovery threshold $r = 0.05$, an average ACID transaction is blocked by a
dirty resource with a probability of 65\%. In addition,
Equation~\ref{eq:blocking-probability} ignores the resources claimed by BASE
parent transactions, which make the probability even higher.

This paragraph shows the additional cost that this implementation of the
adaptive oracle introduces. ACID transactions now have to wait not only on the
locks of other ACID transactions, but also on the resources claimed by BASE
parents and on the dirty resources left by BASE bookkeeper transactions.
Unfortunately, this is unavoidable, because otherwise the isolation promise of
ACID would be broken.

\begin{figure}[H]
  \centering
  \begin{tikzpicture}
    \begin{axis}[
        width=0.9\textwidth,
        height=6.4cm,
        xlabel={dirty fraction $d$},
        ylabel={$P_{\text{block}}$},
        xmin=0, xmax=0.12,
        ymin=0, ymax=1.05,
        xtick={0,0.02,0.04,0.06,0.08,0.10,0.12},
        xticklabel style={/pgf/number format/fixed,
            /pgf/number format/precision=2},
        ytick={0,0.2,0.4,0.6,0.8,1.0},
        yticklabel={\pgfmathparse{100*\tick}\pgfmathprintnumber{\pgfmathresult}\%},
        grid=both,
        grid style={draw=TUMDarkGray!20},
        legend pos=south east,
        legend cell align=left,
        legend style={draw=TUMDarkGray!40, font=\small},
        tick label style={font=\small},
        label style={font=\small},
      ]
      \addplot[TUMBlue, very thick, domain=0:0.12, samples=241]
      {1-(1-x)^50};
      \addlegendentry{worst case, $\ell = L = 50$}
      \addplot[TUMAccentOrange, very thick, dashed, domain=0.0005:0.12,
        samples=241]
      {1-(1-x)*(1-(1-x)^50)/(50*x)};
      \addlegendentry{averaged over $\ell \sim \mathcal{U}\{1,\dots,50\}$}
      \addplot[only marks, mark=*, mark size=1.6pt, TUMBlue]
      coordinates {(0.01,0.395) (0.02,0.636) (0.05,0.923) (0.10,0.995)};
      \node[anchor=north west, font=\scriptsize, TUMBlue]
      at (axis cs:0.0125,0.385) {40\%};
      \node[anchor=north west, font=\scriptsize, TUMBlue]
      at (axis cs:0.0225,0.626) {64\%};
      \node[anchor=north west, font=\scriptsize, TUMBlue]
      at (axis cs:0.051,0.918) {92\%};
      \node[anchor=north east, font=\scriptsize, TUMBlue]
      at (axis cs:0.099,0.99) {99.5\%};
      \addplot[only marks, mark=*, mark size=1.6pt, TUMAccentOrange]
      coordinates {(0.01,0.218) (0.02,0.377) (0.05,0.649) (0.10,0.821)};
      \node[anchor=north west, font=\scriptsize, TUMAccentOrange]
      at (axis cs:0.0125,0.208) {22\%};
      \node[anchor=north west, font=\scriptsize, TUMAccentOrange]
      at (axis cs:0.0225,0.367) {38\%};
      \node[anchor=north west, font=\scriptsize, TUMAccentOrange]
      at (axis cs:0.0525,0.639) {65\%};
      \node[anchor=north west, font=\scriptsize, TUMAccentOrange]
      at (axis cs:0.1025,0.811) {82\%};
    \end{axis}
  \end{tikzpicture}
  \captionsetup{format=plain,justification=centering}
  \caption{Blocking probability against the dirty fraction.}
  \label{fig:blocking-probability}
\end{figure}

\paragraph{Dirty fraction $d_k$ weighted average as one of considered
  designs.}\label{sec:approach-nofilter} One design decision taken along the way
was to smooth the inconsistency level before the oracle looks at it. The
thresholds were not compared against inconsistency level $d_k$ directly, but against an
exponentially weighted moving average (EWMA) $\bar{d}_k$ of it. At every reading
the average was updated with the smoothing factor $\alpha$.

\begin{equation}
  \bar{d}_k \;=\; \alpha \, d_k + (1 - \alpha) \, \bar{d}_{k-1},
  \qquad 0 \le \alpha \le 1, \qquad \bar{d}_0 = d_0
  \label{eq:ewma}
\end{equation}

At $\alpha = 1$ the average is equal to the live value $d_k$, so there is no
smoothing. The baseline smoothing factor was $\alpha = 0.2$. This means that a
new reading $d_k$ contributed only 20\% to the average.

A filtered signal was expected to reduce noise and to avoid rapid mode switches.
That decision was reversed, and the adaptive oracle now reads the live dirty
fraction $d_k$ at the moment of admission. The smoothing was removed for three
reasons.

\begin{bullets}
  \item[Incosistency level $d_k$ is an exact value.] The dirty fraction is an exact number
  maintained by the simulator, and it has no measurement noise. The filter
  therefore had nothing to remove. It only delayed real changes of $d_k$.
  Additionaly, the
  anti-flap work was already done by the hysteresis band
  (Section~\ref{sec:approach-adaptive-oracle}).

  \item[The switch at critical dirty threshold $c$ was unaffected.] The measured $p_{95}$ of the dirty
  fraction stayed within about one percentage point of $c$ at every setting of
  the smoothing factor. The BASE $\to$ ACID switch therefore behaved the same
  way with and without the filter.

  \item[It increased the undershoot of $r$.] The smoothing added a second lag
  to the ACID $\to$ BASE switch, on top of the one described in
  Section~\ref{sec:approach-sampling}. ACID transactions are long, so the time
  between two readings was already large, and with smoothing the switch to BASE
  took even longer. As a result, the system became much cleaner than configured
  before BASE resumed, and the recovery threshold $r$ had little effect.
\end{bullets}


\paragraph{The hybrid mode is designed to work with any
  oracle.}\label{sec:approach-other-oracles} The two oracles of this section are
only two possible policies. The hybrid mode itself does not depend on how a tag
is issued. The dispatchers, the arbiter and the tagging rules work with any
component that implements the \texttt{ModeOracle} interface. A new policy
therefore only needs one new class, which is connected through the hybrid
runtime assembler.

For example, other policies could focus on the throughput or on the length of
the pending queue, or they could use a learned policy like the one used by
Hybrid Transactional Replication~\cite{kobus2018htr}. Every such oracle
also keeps the safety properties of Section~\ref{sec:approach-deadlock}, because
they come from the arbiter and not from the choice of the tag.

\section{Why the hybrid mode is safe}\label{sec:approach-deadlock}

RQ1 asks how a system can add a hybrid mode that adapts between ACID and
BASE \emph{safely}. Safety is not a standard term, so this section first
says what it means in this work and then argues that the design of
Section~\ref{sec:approach-arbiter} has it. The argument is made by
construction. What that buys and what it does not is stated at the end.

\paragraph{What safe means in this
  work.}\label{sec:approach-safety-what} Three properties are meant by the
word. The first one is about the state of a resource. The other two are
about the progress of the run.

\begin{bullets}
  \item[S1. Every resource has one clear state.] A resource is never under
  an ACID lock and marked dirty by BASE at the same time. This is the
  invariant of Section~\ref{sec:approach-arbiter-rule}. It is the
  property that keeps the ACID promise intact, because a lock says
  that the value is exclusive and current, while a dirty mark says
  that an update to that value has been accepted and not yet applied.
  A resource in both states at once would let an ACID transaction
  commit against a value it was promised was the latest, which is
  exactly what it was tagged ACID to avoid
  (Section~\ref{sec:approach-mixing-hazard}).

  \item[S2. Every run is free of deadlock.] The simulator never reaches a
  state in which work is present and no work can proceed. A run
  therefore never stalls, and every simulated step moves the system
  forward.

  \item[S3. No work is thrown away.] Every transaction that is made to wait
  is looked at again and gets a real chance to run. Nothing that is
  parked or deferred is forgotten and nothing is silently dropped.
\end{bullets}

Two things are on purpose not part of this definition. Safety here says
nothing about how long a wait lasts, so a slow run is still a safe run and
the cost of waiting is treated separately in
Section~\ref{sec:approach-blocking-cost}. It says nothing about fairness
either, so one piece of work can be delayed much longer than another, which
is case C2 of Section~\ref{sec:approach-conservatism}. What S3
promises is that the work is still there and is still retried, not that it
runs soon.

\paragraph{S1 holds because neither mode writes without looking
  first.}\label{sec:approach-safety-defined} In a pure run each mechanism
reads only its own record, and that is enough because the other record is
always empty. In a hybrid run both records are in use at once, so each side
has to be able to read the other. The arbiter makes that possible by keying
all three records off the same resource name
(Section~\ref{sec:approach-arbiter-lookup}) and by placing one check on each
of the three paths that write to a record
(Section~\ref{sec:approach-enforcement}).

Those three checks are what give S1. An ACID transaction reads the dirty set
and the claim table before it takes any lock, so it never locks a resource
that BASE is working on. A BASE parent reads the lock table at admission and
takes either all of its claims or none of them, so it never writes over a
resource that ACID holds. A BASE repair reads the lock table before it would
mark its target dirty, so a locked resource is never marked. Every path that
could produce the forbidden state tests for that state first, so the state
has no way of being created.

It is worth being clear about how strong this is. The forbidden state is not
detected and repaired after the fact. It cannot be built at all. That also
means a transaction of either mode keeps what its own tag promises. An ACID
transaction gets exclusive access to its resources even while BASE traffic
runs, which is the reason it was tagged ACID, and a BASE transaction still
commits without waiting for its own repairs, which is the reason it was
tagged BASE. What the hybrid mode changes for a BASE transaction is only
when it is admitted, not what it promises after that
(Section~\ref{sec:approach-base-cost}).

\paragraph{S2 needs a cycle, and only one of the three cases has
  one.}\label{sec:approach-safety-liveness} A deadlock is a cycle in the
wait relation in which every party is waiting while it holds something that
the next party needs. Two conditions therefore have to be met at once. There
has to be a cycle, and each party on it has to hold something.
Figure~\ref{fig:deadlock-freedom} shows the three ways one transaction can
end up waiting for another in a hybrid run, and only the first of them has a
cycle. That first one comes from the base simulator, which already breaks it
by aborting one group. The two kinds of waiting added by this thesis, the
waiter registry and the two deferred queues, both park work that holds
nothing, so neither of them can be part of a cycle.

\begin{figure}[htbp]
  \centering
  \includegraphics[width=\textwidth]{fig-deadlock-freedom.pdf}
  \captionsetup{format=plain,justification=centering}
  \caption{The three kinds of waiting in a hybrid run.}
  \label{fig:deadlock-freedom}
  \medskip
  \begin{figurekeys}
    \item[1] The cycle already exists in the received simulator. Cautious
    waiting detects it and aborts one of the two groups, which releases
    every lock that group holds.
    \item[2] The ACID transaction is parked as a waiter before it is admitted
    and before it takes any lock. No other transaction can wait behind
    it.
    \item[3] A deferred entry holds nothing and joins no lock queue. It moves
    on only once every resource it needs is free again.
  \end{figurekeys}
\end{figure}

\paragraph{ACID waiting for ACID is not new.} The first case is two ACID
groups that each need a lock the other one holds. This is the situation of
Section~\ref{sec:approach-group-locking}, and it is present in the base
simulator before anything in this thesis is added. A group takes its locks
in two stages, so a cycle between groups can form. The base simulator does
not prevent it. It detects it and aborts one of the two groups, which
releases every lock that group held. This thesis changes nothing about that
mechanism, so this case is neither better nor worse than it was.

\paragraph{ACID waiting for BASE work runs in one direction.} The second
case is the new one. An ACID transaction that meets a dirty or claimed
resource goes into the waiter registry and waits until the BASE work on that
resource is done. This edge cannot lie on a cycle, for two reasons.

The first is that the edge has no partner in the other direction. A BASE
transaction never waits on an ACID transaction. When BASE meets an ACID lock
it backs off instead, which is the third case below. So there is no edge
from the BASE side back to the ACID side that could close the loop.

The second is that the chain ends. The thing the ACID transaction is waiting
for is a repair or a BASE parent, and both of those need only a reconciler
processor or an execution processor to finish. Neither of them needs
anything that an ACID transaction could be holding. The last node in the
chain therefore waits for nobody, and a chain that ends in a node that waits
for nobody is not a cycle.

One detail is worth naming, because it is the one case where the waiting
ACID side is not empty handed. A fresh root transaction is tested at
admission and holds nothing at all. A bookkeeper is tested later, when it
takes its own locks, and at that moment its group may already hold locks
from the root. Another ACID group can then wait behind those locks. That
still does not create a cycle, because the chain from the parked group runs
on into a reconciler, and the argument above applies unchanged at that
point. It only makes the chain one step longer.

\paragraph{BASE meeting an ACID lock is a back off, not a wait.} The third
case is a BASE parent or a BASE repair whose resource is ACID locked. It is
placed on a deferred queue, and the important part is what it is holding
while it sits there. A parent takes all of its claims or none of them, so a
deferred parent has claimed nothing. A deferred repair has marked nothing
dirty. Neither of them joins a lock queue. Nothing they need is held back
from anyone else, so no other transaction can end up waiting behind a
deferred entry. A node that nothing points to cannot be a link in a cycle.

\paragraph{The three cases together.} Adding the arbiter therefore adds
waiting but no new cycle, which is S2. Two further points follow from that.
The number of aborts does not change, because the two abort rules of
Section~\ref{sec:approach-group-locking} read the lock table only, and a
transaction parked in the waiter registry holds no lock that those rules
would look at. And the cost of the new waiting shows up as time rather than
as wasted work, which is what Section~\ref{sec:approach-blocking-cost}
measures.

\paragraph{S3 holds because every waiting entry is looked at
  again.}\label{sec:approach-safety-progress} Freedom from deadlock is not
enough on its own. A run could also lose work if a waiting entry were parked
and never looked at again. The two waiting structures avoid that in two
different ways, and the difference is deliberate.

The waiter registry is woken. Two events free a resource for ACID, a repair
finishing and a BASE parent completing, and both of them notify the
registry, which then retries every transaction parked on that resource.
Nothing polls here, because the two events that matter are both visible at
the moment they happen.

The deferred queues are polled, because the event that matters to them is
not visible in the same way. Nothing tells a deferred entry that a lock has
gone away, so each queue is swept and every entry is asked the same question
again. An entry that is still blocked goes to the back of the queue and is
asked again on the next sweep, so no entry is tested only once. The sweep
that runs once per simulated tick is the one that makes this complete. A
lock is also released when a group aborts, and an abort is not a completion,
so a sweep tied only to completions could leave an entry deferred after its
reason for being deferred was already gone. Sweeping every tick removes that
dependency.

Nothing is dropped on either path. A deferred repair is created in full when
its turn comes, with the one exception that a sibling repair may have
dirtied the same resource in the meantime, in which case the deferred one is
collapsed into it. That is the same fan-out collapse a repair meets in a
pure BASE run, so it removes duplicated work and not the work itself.
Together these rules give S3. Whatever the reason a piece of work was made
to wait, that work is either woken by the event that removed the reason or
retested within one tick.

\paragraph{What this argument does and does not
  establish.}\label{sec:approach-safety-limits} The three properties are
established by construction and not by measurement. S1 holds because the
three checks sit on the three write paths and there is no fourth path. S2
and S3 hold because of the shape of the wait relation those checks produce.
No runtime assertion checks any of this during an experiment, so the claim
to make is that the unsafe state cannot be built, and not that it was looked
for and never seen. The run data does record how often each of the two
checks fired, and Section~\ref{sec:future-work} names reporting those counts
as the natural way to show that the mixed state was reached and resolved
constantly rather than never reached at all.

The two known imprecisions of Section~\ref{sec:approach-conservatism} do not
weaken the three properties, because both of them err in the same direction.
Under C1 coarse lock granularity makes the BASE checks look at a whole
bucket of resources instead of one, so work can be deferred that did not
have to be, which costs time and never breaks S1. Under C2 a deferred entry
whose resource stays locked for a long time can be delayed for a long time.
That is a starvation case rather than a deadlock, so S2 and S3 still hold.
